% ============================================================================
% supplementary/supp.tex — ICME 2027 supplementary material (Plan B version, 2026-11-06)
% ----------------------------------------------------------------------------
% Rules (official ICME 2026 kit, re-verified 2026-11-06): supplemental material
% must be ANONYMIZED, zipped into ONE file; the kit states a 20 MB limit while the
% website's submission list says 50 MB -> cite both, treat 20 MB as binding until
% the 2027 rules are published. Reviewers are not obliged to read it.
% Compile: pdflatex supp.tex (standalone; NOT \input into main.tex).
%
% Plan B status (2026-11-06): the paper is repositioned as a development benchmark
% + diagnostic study + a controlled-intervention protocol. S1-S3 now carry their
% final numbers; S4-S8 carry the material their results support. The old-script
% +0.50 dB comparison and its per-scene figure are NOT results: they appear only
% in S3 as a single historical observation.
%
% Main-text migrations (Lead-approved 2026-11-06):
%   * Table I (layer asymmetry)   -> S2 Table S1
%   * unified control full grid   -> S3 Tables S2-S3 (main text keeps one compact table)
%   * old Fig.3 (per-scene dPSNR of the old-script comparison) -> WITHDRAWN
% ============================================================================
\documentclass[10pt]{article}
\usepackage[margin=1in]{geometry}
\usepackage{amsmath,amssymb}
\usepackage{graphicx}
\usepackage{booktabs}
\usepackage{array}
\usepackage[hidelinks]{hyperref}
\newcommand{\Om}{\Omega}
\newcommand{\dBA}{\,\mathrm{dB}}
\newcommand{\quiet}{\mathcal{L}_{\mathrm{quiet}}}
\title{Supplementary Material for\\
What Filmed Glass Does to Single-Image Reflection Removal:\\
A Development Benchmark and a Diagnostic Study}
\author{Anonymous ICME submission}
\date{}
\begin{document}
\maketitle

\noindent This document accompanies the paper. S1 documents the data and capture
protocol; S2 the imaging and diagnostic measurements; S3 the unified-recipe control
experiment, its full table and its run-level reproducibility; S4 the negative results;
S5 the supervision placement and mask-quality study; S6 the light-weight-method
motivation; S7 the frozen independent-capture protocol; S8 the reproducibility and
release boundary. All numbers are reported with the pool and protocol that produced
them.

% ============================================================================
\section{Dataset statistics and protocol}\label{supp:data}
% All numbers carry: pool name / #pairs / #scenes / QC rule.
\subsection{Pools}
\begin{itemize}
\item Training pools: the real branch is drawn from sessions with paired
  reflection/black-cloth captures; the synthetic branch follows the public SIRR
  mixture; the public branch uses released paired data. The exact per-pool counts are
  released with the training code.
\item The retired June-44 subset (two scenes) is used only in the disclosed
  configuration; its replacement pool is the rebuilt one, and the main configuration's
  composition is disclosed (Lead decision D3-a: disclosure-only).
\item Development panel (DP-15): 15 scenes, $1280{\times}832$, one ground-truth frame
  per scene, median 374 frames between an input and its ground truth (6\% within 50),
  hand-annotated occluder region $\Om$ (median 4.6\%, range 2.9--9.2\%, fixed
  intensity rule, tree-shade blind spot stated in the main text).
\end{itemize}
\subsection{Capture defects and the QC protocol}
\begin{itemize}
\item Exposure lock: with auto-exposure active, $\mathrm{median}(I/T)=0.99$ and the
  residual layer correlates with image brightness at $r=0.32$; after locking,
  $r=-0.01$.
\item Alignment: on one batch $88.7\%$ of the pairs were shifted by more than 3\,px
  (median 8.2\,px). Per-pair translation alignment is a single-seed screening result
  on that pool and is not used for any main-table claim.
\item Privacy: faces and licence plates are redacted before release; the release is
  research-only.
\end{itemize}

% ============================================================================
\section{Imaging and diagnostic details}\label{supp:diag}
\subsection{Layer statistics on four datasets (moved Table I)}
\begin{table}[h]\centering
\caption{Layer statistics of the transmitted ($T$) and reflected ($R$) layers on
Real20, SIR$^2$-Obj, SIR$^2$-Wild and the filmed-glass panel. The filmed column is
measured identically; the public-dataset columns report the observed range across the
three datasets (per-dataset values in the released measurement script).}
\begin{tabular}{@{}lccc@{}}
\toprule
Quantity & public range & Real20--SIR$^2$ & Ours (filmed) \\
\midrule
sharp($T$) & 0.041--0.062 & -- & \textbf{0.027} \\
sharp($R$)/sharp($T$) & 2.51--3.32 & -- & \textbf{6.76} \\
chroma of $T$ & 0.322--0.393 & -- & \textbf{0.219} \\
mean brightness & ${}\approx{}$Real20 & -- & ${}\approx{}$Real20 \\
mean $|R|$ & ${}\approx{}$Real20 & -- & ${}\approx{}$Real20 \\
HF fraction of $R$ & ${}\approx{}$Real20 & -- & ${}\approx{}$Real20 \\
\bottomrule
\end{tabular}
\end{table}

\subsection{Three-band decomposition inside/around $\Om$}
\begin{table}[h]\centering
\caption{Three-band decomposition on the development panel. Band values are not
comparable across bands (PSNR depends on band brightness); only the within-band
contrast against identity is informative.}
\begin{tabular}{@{}lcccc@{}}
\toprule
Band & Area & $\mathbb{E}|I-T|$ & Identity PSNR (dB) \\
\midrule
$\Om$ (hand) & \textbf{4.6\%} (2.9--9.2) & 10.3 & 24.52 \\
annulus (15\,px) & 44.4\% & 44.4 & 13.76 \\
far field & remainder & 16.8 & 19.52 \\
\bottomrule
\end{tabular}
\end{table}

\subsection{Structural damage and the content-agnostic ceiling}
\begin{itemize}
\item Structure retained relative to ground truth after a per-image affine
  correction, inside $\Om$: 1.14 (identity), 0.50 (fine-tuned baseline), 0.65
  (regularizer arm); affine-corrected MAE 10.51, 12.82 and 11.70 gray levels (these
  are single-seed development measurements).
\item Content-agnostic ceiling: a Gaussian smoothing of the input is the best a
  content-blind method can do. On the weakly reflected 59\% of $\Om$ the ceiling is
  29.50\,dB against identity's 31.45\,dB; on the strongly reflected 7\% it is
  15.24\,dB against 14.46\,dB. No trained model beats a plain blur in either regime.
\end{itemize}

\subsection{Frequency overlap and the input-HF splicing test}
\begin{itemize}
\item 34.5\% of the reflection energy lies above a $\sigma{=}8$ low-pass, so the
  reflection is not a veil that can be filtered away.
\item Replacing the output's high frequencies by those of the input restores
  structure inside $\Om$ but lowers PSNR: 22.26$\rightarrow$22.04$\rightarrow$21.66\,dB
  at $\sigma{=}2,4$ (regularized arm), and 23.73$\rightarrow$23.39$\rightarrow$22.97\,dB
  for a 50:50 mixture.
\end{itemize}

\subsection{Per-image photometric oracle}
A per-image affine oracle fitted on each image's ground truth recovers
0.15--0.26\,dB, whereas a single shared affine fitted on the other 14
development-panel scenes (leave-one-scene-out) recovers 0.07\,dB. Fine-tuning
therefore consumes the transferable part of the photometric error and leaves a
residual that is per-image; any module that estimates a global photometric quantity
from a single frame is bounded by the shared value.

% ============================================================================
\section{The unified-recipe control experiment}\label{supp:control}
\subsection{Protocol}
\begin{itemize}
\item One recipe for all arms: same initialization and seeds (101, 202, 303), same
  data mixture (30\% real / 40\% synthetic / 30\% public), 6000 iterations, learning
  rate $10^{-4}$, crop curriculum 224$\rightarrow$256$\rightarrow$320; only the
  regularizer term differs.
\item Arms: base (no extra term); identity ($\lambda\|\widehat T-I\|_1$,
  $\lambda=0.409$); identity-m (same, $\lambda=0.085$, magnitude-matched to the
  change-aware weight); hard (change-aware statistic thresholded at $\tau=0.05$);
  shuffled (change-aware weights permuted across pixels); full ($\quiet$,
  $\lambda=1$, $\tau=0.05$).
\item Evaluation: frozen uint8 script on the 15-scene development panel; scene-level
  paired differences with 95\% CI over the 15 scenes.
\end{itemize}

\subsection{Full control table (moved from the main text)}
\begin{table}[h]\centering
\caption{Unified-recipe arms on the development panel (three seeds). $\Delta$ is the
scene-level paired mean against \emph{full}; the last column counts scenes where
\emph{full} is higher.}
\begin{tabular}{@{}lcccc@{}}
\toprule
Arm & PSNR (dB) & $\Delta$ vs full (dB) & 95\% CI & scenes \\
\midrule
base & $22.38\pm0.23$ & $+0.19$ & $[+0.05,+0.33]$ & 11/15 \\
identity ($\lambda{=}.409$) & $22.49\pm0.34$ & $+0.079$ & $[+0.021,+0.137]$ & 14/15 \\
identity-m ($\lambda{=}.085$) & $22.59\pm0.21$ & $-0.014$ & $[-0.215,+0.186]$ & -- \\
hard ($\tau{=}.05$) & $22.82\pm0.36$ & $-0.246$ & $[-0.355,-0.138]$ & 2/15 \\
shuffled & $22.36\pm0.39$ & $+0.209$ & $[+0.033,+0.386]$ & -- \\
full ($\lambda{=}1$, $\tau{=}.05$) & $22.57\pm0.00$ & -- & -- & -- \\
\bottomrule
\end{tabular}
\end{table}

Two entries in this table look like copying or indexing errors, but both are genuine
averaging coincidences. The three \emph{full}-arm checkpoints were trained
independently and have pairwise different SHA-256 fingerprints; their per-seed
scores are 22.5698, 22.5786 and 22.5721\,dB, which is why the table reports
$\pm0.00$ (spread $\approx$0.004\,dB). Likewise, the hard-arm seeds 101 and 202
both average 23.0268\,dB although their per-scene PSNR differs by 0.02--1.12\,dB
and cancels in the mean; their checkpoints and prediction files also differ
pairwise. A small seed spread is therefore not evidence of distinguishability:
the main text reads arm differences against the 0.88\,dB run-level spread (below),
not against the $\pm0.00$ of a single arm.

\subsection{Run-level reproducibility}
\begin{itemize}
\item Same arm, same seed, same script, re-run: base $22.998\rightarrow22.119$
  ($-0.88$\,dB); full $23.079\rightarrow22.570$ ($-0.51$\,dB).
\item Cross-date contrasts: $+1.008$\,dB (base) and $+0.369$\,dB (full).
\item Consequence: run-level uncertainty is at least as large as seed-level
  uncertainty; differences of $\sim$0.08\,dB cannot be attributed to a regularizer,
  and the main text uses this bound when interpreting the control.
\end{itemize}

\subsection{Historical observation from the pre-unified scripts}
Before the recipes were unified, an earlier comparison of the two training scripts
produced $+0.50$\,dB. Because the scripts differ in recipe and match, this number is a
historical observation, not a result, and it is not comparable with the unified table;
the per-scene figure built from it is withdrawn. Paired differences between
LightRR's re-evaluation and the old three-seed references of those scripts
(22.99 / 22.49) are likewise historical and appear only in S6.

\subsection{What the control supports}
The control bounds what this family of training-time regularizers achieves on the
development panel (a few tenths of a decibel, within the run-level spread) and shows
that the specific change-aware shape is not the lever: identity-m is on par with
full, and hard is significantly better. It supports the benchmark's ability to make
this comparison controlled; it does not support a performance claim for any arm.

% ============================================================================
\section{Negative results}\label{supp:negative}
\begin{itemize}
\item Perceptual-loss weighting did not close the structural gap; the weighting
  shifted the trade-off but not the total.
\item Splicing input high frequencies into the output restored detail and lowered
  PSNR (S2).
\item Mixing occluder-region data raised region-restricted PSNR while degrading the
  residual-reflection coefficient.
\item Photometric adapters and residual gates were rejected by cheap probes before
  full training.
\item Averaging three independently trained models failed because the seeds converge
  to different basins.
\end{itemize}

% ============================================================================
\section{Supervision placement and mask quality}\label{supp:placement}
\begin{itemize}
\item 82-pair five-region placement study (three seeds unless noted): no region
  17.76 / hand-$\Om$ 19.41 / random disk 19.90 / fixed disk 20.25 / whole image 20.05.
  The hand annotation does not beat an equal-area random disk; occluder-region
  specificity is therefore withdrawn.
\item Mask quality: hand annotation vs a supervision-only detector IoU 0.376—0.063;
  no downstream gain from a better detector, consistent with the region-specificity
  withdrawal.
\item $\Om$ annotation rule and blind spot (tree-shade pixels at threshold 15); it is
  an annotation of where the reflection path was blocked, not a physical ground truth.
\end{itemize}

% ============================================================================
\section{Light-weight-method protocol and results}\label{supp:lightrr}
\begin{itemize}
\item Protocol: LightRR is fine-tuned in-domain with its official recipe for 102,500
  iterations (41 epochs; one epoch more than pre-registered, disclosed) on the same
  training pools and evaluated with the same frozen script on the development panel.
\item Two learning-rate runs: $10^{-4}\rightarrow20.221$\,dB,
  $5\times10^{-5}\rightarrow20.198$\,dB; best-of-two $20.221$\,dB; an earlier
  official-recipe rerun scored 19.175\,dB.
\item Context: more than 2\,dB below the strongly fine-tuned baseline, winning on no
  scene (0/15); the residual-reflection coefficient moves from $-0.859$ zero-shot to
  only $-0.489$/$-0.511$, against $-0.29$ to $-0.34$ for the fine-tuned RDNet.
\item Historical paired differences against the old three-seed references
  (22.99\,dB and 22.49\,dB): $-2.766$\,dB (95\% CI $[-3.371,-2.162]$) and
  $-2.264$\,dB; these are historical and are not results of the unified protocol.
  Historical reference only; the unified-protocol paired $\Delta$ will be reported in
  the camera-ready version if available.
\end{itemize}

% ============================================================================
\section{Frozen protocol for an independent capture}\label{supp:indep}
\begin{itemize}
\item New A/B capture through the same film: 6--8 indoor and 4--6 outdoor scenes at
  sites disjoint from the development sessions, same camera and film, locked
  exposure/white balance, black-cloth ground truth, QC $|g-1|\le0.15$ and
  phase-correlation shift $\le2$\,px.
\item Frozen before scoring: model checkpoints, baselines, evaluation script, primary
  contrast and the scene-level paired CI procedure. No tuning on the held-out set.
\item One-shot evaluation; per-scene results reported; if the primary contrast
  crosses zero, the claim shrinks instead of being re-tuned (the pre-registered
  analysis plan is included in the release).
\end{itemize}

% ============================================================================
\section{Reproducibility and release boundary}\label{supp:repro}
\begin{itemize}
\item Anonymised artifact statement: code, evaluation script, split hashes and the
  frozen analysis plans are releasable; training data contains faces and plates and
  is released only after redaction.
\item Correction log for the 2026-10-29 release package: retracted numbers are listed
  in the release's correction file; the old-script $+0.50$\,dB comparison and its
  figure are withdrawn from this paper. No external links appear in the
  double-blind copy.
\end{itemize}

\end{document}
